\section{The joint cost of dimension, horizon, and confidence}
\label{sec:confidencelearning80}

The fixed-confidence interior law has a sharp extension to arbitrary
confidence.  The dependence on dimension and on the confidence
logarithm is additive after their respective dimensional factors are
included.  In this section the unknown measurement is still ordered,
binary, consuming, and classical-output.  Put
$\mathfrak I_d=\{E:I_d/4\preceq E\preceq3I_d/4\}$ as in
Section~\ref{sec:interiorcoding79}.  Logarithms in call bounds are
natural unless indicated otherwise.

\begin{theorem}[Uniform interior learning at every confidence]
\label{thm:interiorconfidence80}
There are absolute $c,C>0$ such that, for every $d,N\ge1$,
$0<\delta\le2^{-13}$, and $0<\eta\le1/8$, the minimax
every-record call budget for
\begin{equation}\label{eq:interiorconfidencerisk80}
 \sup_{E\in\mathfrak I_d}
 \mathbb P_E\{d_N(\mathcal M_E,\mathcal M_{\widehat E})>\delta\}
 \le\eta
\end{equation}
is between
\begin{equation}\label{eq:interiorconfidence80}
 c\,\frac{N}{\delta^2}
       \left(d^2+d\log\frac1\eta\right)
 \quad\hbox{and}\quad
 C\,\frac{N}{\delta^2}
       \left(d^2+d\log\frac1\eta\right).
\end{equation}
The lower bound permits arbitrary coherent adaptive training,
references, quantum memory, and public bounded stopping.  The estimate
may be any legal effect.  The upper bound uses independent one-call
Choi preparations and collective processing of completed outputs.
Consequently the same order, with absolute constants, holds for
every fresh-block cap $1\le b\le N$.
\end{theorem}

The proof combines the dimension lower bound of
Theorem~\ref{thm:dimensionlower78} with a separate confidence
obstruction and an amplification of the estimator in \cite{MB67}.
The latter procedure is the tomography primitive throughout; the
amplification uses its dimension-dependent tail guarantee.

\begin{lemma}[Confidence requires observations in every coordinate]
\label{lem:diagonalconfidence80}
In the parameter range of Theorem~\ref{thm:interiorconfidence80},
every learner satisfying \eqref{eq:interiorconfidencerisk80} has
\begin{equation}\label{eq:diagonalconfidence80}
 M\ge\frac{dN}{6\cdot512^2\delta^2}\log\frac1\eta.
\end{equation}
This bound already holds when the unknown effects are diagonal in
a fixed known basis.
\end{lemma}
\begin{proof}
Set $\Delta=512\delta/\sqrt N\le1/16$ and consider
\[
 E_0=I/2,\qquad E_i=I/2+\Delta|i\rangle\langle i|,
 \quad 1\le i\le d.
\]
All these effects belong to $\mathfrak I_d$.  For each pair
$(E_0,E_i)$, repeating the input $|i\rangle$ gives Bernoulli
parameters $1/2$ and $1/2+\Delta$.  The comparison quantity in
Lemma~\ref{lem:bernoulliproduct75} satisfies
\[
 T_N\ge\tfrac12\sqrt N\,\Delta=256\delta.
\]
That lemma, together with $256\delta\le1$, therefore gives
\begin{equation}\label{eq:diagonalseparation80}
 d_N(\mathcal M_{E_0},\mathcal M_{E_i})\ge4\delta.
\end{equation}

We bound the information used to distinguish all these alternatives
from the same null effect.  For a diagonal effect with entries
$p_1,\ldots,p_d$, the reference-assisted output is unchanged if
the input is first measured in the known basis, yielding $J=j$,
and a Bernoulli bit of parameter $p_j$ is then drawn.  We retain
the basis label in an extra register for the analysis.  Discarding
it gives the original device, so any original learner, including
its final estimate, is recovered by ignoring these extra registers.
This is the diagonal classicalisation used in
\cite[Lemma~IV.19]{MB67}; the following relative-entropy calculation
also treats different quantum inputs under the two hypotheses.

For an input-reference state $\rho_{AR}$, write
$\tau_j=\langle j|\rho_{AR}|j\rangle$.  The revealed output is
\[
 \widetilde{\mathcal M}_{p}(\rho_{AR})
 =\sum_{j=1}^d |j\rangle\langle j|\otimes
      \bigl((1-p_j)|0\rangle\langle0|
                 +p_j|1\rangle\langle1|\bigr)\otimes\tau_j.
\]
Let $\rho_0,\rho_i$ be arbitrary current input-reference states,
where the reference includes all retained memory.  The common
channel $\rho\mapsto\sum_j|j\rangle\langle j|\otimes\tau_j$
cannot increase relative entropy.  Appending the coin conditional
on $j$ adds its classical relative entropy with weight
$\operatorname{tr}\tau_j^0$.  Thus the classical--quantum
relative-entropy identity gives
\begin{equation}\label{eq:diagonalchain80}
 \begin{split}
 D\bigl(\widetilde{\mathcal M}_{0}(\rho_0)
       \,\big\|\,\widetilde{\mathcal M}_{i}(\rho_i)\bigr)
 &\le D(\rho_0\|\rho_i)
       +\kappa_\Delta\operatorname{tr}\tau_i^0,\\
 \kappa_\Delta
 &=D\bigl(\operatorname{Bern}(1/2)
            \,\big\|\,\operatorname{Bern}(1/2+\Delta)\bigr)\\
 &=-\tfrac12\log(1-4\Delta^2)\le3\Delta^2 .
 \end{split}
\end{equation}
The subscript on $\widetilde{\mathcal M}$ now denotes the
hypothesis.  This calculation uses no equality or independence
assumption for $\rho_0,\rho_i$.  Starting from the common initial
state, it also establishes finite relative entropy at every stage.

Keep each revealed label without feeding it into the original
controls.  All intervening operations, including classically
controlled or coherent ones, are common channels.  Pad a stopped
record to its every-record cap $M$ using ignored dummy calls.
Let $T_i$ count the revealed labels equal to $i$ among these $M$
calls, and let $\Omega_0,\Omega_i$ denote the complete final states
with the extra labels retained.  Iterating
\eqref{eq:diagonalchain80} yields
\[
 D(\Omega_0\|\Omega_i)\le3\Delta^2\mathbb E_0T_i,
 \qquad
 \sum_{i=1}^d\mathbb E_0T_i=M.
\]
The same null experiment appears for every $i$, which is essential
for summing these inequalities.  General measurable records are
included by the relative-entropy chain rule for classical--quantum
states; no countability assumption on the learner's observations
is required.

By \eqref{eq:diagonalseparation80}, choose hypothesis $0$ precisely
when the returned effect lies in its closed radius-$\delta$ ball.
This is a common binary test with both errors $\alpha_i,\beta_i$
at most $\eta$.  The test is measurable by the hybrid continuity
bound for $d_N$.  Data processing to this test gives
\[
 \begin{split}
 D(\Omega_0\|\Omega_i)
 &\ge D\bigl(\operatorname{Bern}(1-\alpha_i)
                  \,\big\|\,\operatorname{Bern}(\beta_i)\bigr)\\
 &\ge(1-\eta)\log(1/\eta)-\log2
 \ge\tfrac12\log(1/\eta),
 \end{split}
\]
where $\eta\le1/8$ was used.  Summing over $i$ and substituting
$\Delta$ proves \eqref{eq:diagonalconfidence80}.
\end{proof}

\begin{lemma}[Amplification using a dimension-dependent seed]
\label{lem:binaryconfidenceupper80}
For all $d\ge1$, $0<\epsilon\le1$, and $0<\eta\le1/8$,
there is a common estimator on $\mathfrak E_d$ with legal output
$F$ such that
\begin{equation}\label{eq:binaryconfidenceupper80}
 \sup_{E\in\mathfrak E_d}
       \mathbb P_E\{\|F-E\|_{\rm op}>\epsilon\}\le\eta,\qquad
 M\le C\epsilon^{-2}
             \left(d^2+d\log\frac1\eta\right).
\end{equation}
The acquisition uses independent one-call input/reference pairs.
The output is a Borel function of collective measurements on their
completed outputs.
\end{lemma}
\begin{proof}
Corollary~III.9 of \cite{MB67} gives a legal binary-effect estimator
with operator error $u$ and failure $q$ using
$O(u^{-2}(d^2+d\log(1/q)))$ calls when
$q>4e^{-4d^2}$.  Choose the rational seed probability
\[
 q_d=2^{-3d}.
\]
It satisfies that strict condition for every $d\ge1$, since
$4d^2-3d\log2-\log4>0$.  A seed run with $u=\epsilon/3$
therefore costs at most $Cd^2\epsilon^{-2}$ calls.

Let $k$ be the smallest positive odd integer with
$k\ge2d^{-1}\log_2(1/\eta)$.  Make $k$ independent seed runs,
obtaining legal effects $F_1,\ldots,F_k$.  Select the first index
$j$ for which more than $k/2$ of the effects have operator distance
at most $2\epsilon/3$ from $F_j$.  If there is no such index,
return $F_1$.  Norm continuity and the fixed tie rule make this
selection Borel, and it always returns a legal effect.

If more than half the seed estimates are within $\epsilon/3$
of $E$, a selectable index exists.  Every selectable index has a
neighbour among these good estimates, and so its distance from
$E$ is at most $\epsilon$.  The probability of at least $k/2$
failed seeds is at most
\[
 2^k q_d^{\,k/2}
 =2^{k-3dk/2}\le2^{-dk/2}\le\eta.
\]
This proves the risk bound.  Since
$k<2d^{-1}\log_2(1/\eta)+2$, the total call count has the order
in \eqref{eq:binaryconfidenceupper80}.  Each seed uses fresh Choi
preparations; all subsequent joint processing acts on completed
outputs.  The amplification is a modification of the repetition
option in \cite[Remark~III.4]{MB67}, with seed failure chosen at
the dimension scale rather than fixed independently of $d$.
\end{proof}

\begin{proof}[Proof of Theorem~\ref{thm:interiorconfidence80}]
Theorem~\ref{thm:dimensionlower78} gives
$M\ge c_1d^2N\delta^{-2}$ throughout the stated range.
Lemma~\ref{lem:diagonalconfidence80} supplies
$M\ge c_2dN\delta^{-2}\log(1/\eta)$.
The maximum of these two bounds is at least half their sum after
decreasing the absolute constant, which proves the lower bound in
\eqref{eq:interiorconfidence80}.

For the upper bound apply Lemma~\ref{lem:binaryconfidenceupper80}
at $\epsilon=\delta/(8\sqrt N)$.  On its good event,
$\|F-E\|_{\rm op}\le\epsilon\le1/8$.  Since $E\in\mathfrak I_d$,
Weyl's inequality puts both endpoints in $[I/8,7I/8]$.
Lemma~\ref{lem:interiormetric78} gives
$d_N(E,F)\le4\sqrt N\,\epsilon=\delta/2$.
The call count is the required upper bound.  The procedure is in
the fresh one-call class, whereas the lower bound imposed no
restriction on coherent training.  This also proves the assertion
for every fresh-block cap.
\end{proof}

For $d=1$, \eqref{eq:interiorconfidence80} reduces to
$\Theta(N\delta^{-2}\log(1/\eta))$, the scalar order, because
$\log(1/\eta)\ge\log8$.  For growing $d$, the confidence term
becomes comparable to the dimension term when
$\log(1/\eta)$ is of order $d$.

\begin{corollary}[Binary operator-norm learning at every confidence]
\label{cor:operatorconfidence80}
For $d\ge1$, $0<\epsilon\le2^{-14}$, and $0<\eta\le1/8$,
the minimax every-record call budget for operator-norm error
$\epsilon$ on the complete ordered binary effect body is
\begin{equation}\label{eq:operatorconfidence80}
 \Theta\!\left(\epsilon^{-2}
                   \left(d^2+d\log\frac1\eta\right)\right),
\end{equation}
with absolute constants.  The lower bound permits arbitrary
coherent adaptive training; the upper bound uses independent Choi
preparations and collective completed-output processing.
\end{corollary}
\begin{proof}
The upper bound is Lemma~\ref{lem:binaryconfidenceupper80}.
For the lower bound restrict to $\mathfrak I_d$ and use
Theorem~\ref{thm:interiorconfidence80} with $N=1$ and
$\delta=2\epsilon$, since $d_1(E,F)=2\|E-F\|_{\rm op}$.
\end{proof}

On the complete effect body, the confidence obstruction also
combines with the independent-block obstruction.  For $d\ge2$,
$1\le b\le N$, $0<\delta\le\min\{2^{-13},\delta_*\}$, and
$0<\eta\le1/8$, Theorem~\ref{thm:blocklearning78} and the
interior subfamily give
\begin{equation}\label{eq:fullbodyconfidencelower80}
 M_b^\star(d,N,\delta,\eta)
 \ge c\delta^{-2}
       \left[d^2N+\left(dN+\frac{N^2}{b}\right)\log\frac1\eta\right].
\end{equation}
A second upper construction follows from
Lemma~\ref{lem:binaryconfidenceupper80} at operator accuracy
$\epsilon=\delta/(2N)$.  The hybrid inequality
$d_N(E,F)\le2N\|E-F\|_{\rm op}$ converts its good event to the
required future loss.  It uses only fresh one-call preparations,
so it is admissible for every block cap $b$.  Combining its public
call budget with the construction in
Theorem~\ref{thm:blocklearning78} gives, in the same parameter range,
\begin{equation}\label{eq:fullbodyconfidenceupper80}
 M_b^\star(d,N,\delta,\eta)
 \le C\frac{N^2}{\delta^2}
       \min\left\{d^2+d\log\frac1\eta,\,
                    \frac{d^4}{b}\log\frac d\eta\right\}.
\end{equation}
The learner may choose the construction with the smaller public
every-record budget before making any calls.  Both constructions
are common to all unknown effects.  The two-sided bounds
\eqref{eq:fullbodyconfidencelower80}--\eqref{eq:fullbodyconfidenceupper80}
do not in general match when dimension and horizon grow together.

\begin{corollary}[Jointly optimal learned descriptions at every confidence]
\label{cor:interiorconfidencecode80}
For $d,N\ge1$, rational $0<\delta\le2^{-13}$, and
$0<\eta\le1/8$, one common learner on $\mathfrak I_d$ returns
a fixed-length binary word whose public deterministic decoder is
legal and satisfies \eqref{eq:interiorconfidencerisk80}, with
\begin{equation}\label{eq:interiorconfidencecode80}
 M=O\!\left(\frac{N}{\delta^2}
                  \left(d^2+d\log\frac1\eta\right)\right),
 \qquad
 B=d^2\log_2(\sqrt N/\delta)+O(d^2).
\end{equation}
All constants are absolute.  Acquisition uses independent one-call
Choi preparations, and encoding uses no additional device calls.
Every learner with these risk and decoder conventions requires
the lower call order in \eqref{eq:interiorconfidence80} and
$B\ge d^2\log_2(\sqrt N/\delta)-11d^2$.
\end{corollary}
\begin{proof}
Put $k_N=\lceil\sqrt N\rceil$.  Apply
Lemma~\ref{lem:binaryconfidenceupper80} at
$\epsilon=\delta/(64k_N)$ and failure $\eta$, and clip its legal
output $F$ spectrally to $[I/4,3I/4]$, obtaining $T$.
On the good event, Weyl's inequality gives
$\|T-F\|_{\rm op}\le\epsilon$ and hence
$\|T-E\|_{\rm op}\le2\epsilon$.
This uses only the eigenvalue displacements of $F$ and does not
assume operator-Lipschitz spectral clipping.  The interior
comparison yields $d_N(E,T)\le\delta/8$.

Use the exact dictionary of Theorem~\ref{thm:interiorcodec79}
at the same $(d,N,\delta)$.  As in the proof of
Theorem~\ref{thm:interiorlearnedcode79}, spectral clipping and
the specified coordinate rounding give a Borel finite-valued
index map, which can be incorporated into the final collective
measurement.  The decoded centre satisfies
$d_N(T,C)\le5\delta/16$, so its total error is at most
$7\delta/16<\delta$ on the original good event.
Thus the finite output does not increase the failure probability
or supply an exact-real side channel to the decoder.

The bound $k_N^2\le4N$ and Lemma~\ref{lem:binaryconfidenceupper80}
give the call count.  Theorem~\ref{thm:interiorcodec79} gives the
payload.  Conversely, a decoded word is an estimator, so
Theorem~\ref{thm:interiorconfidence80} applies to its calls.
For each target, positive success probability forces some decoded
effect into its radius-$\delta$ ball.  The fixed decoder's range
is therefore a cover, and Theorem~\ref{thm:interiorentropy79}
gives the stated bit lower bound.
\end{proof}

These upper bounds use the ideal trusted tomography operations of
\cite{MB67}.  The Borel finite readout is a statistical construction;
it alone supplies no gate synthesis or reconstruction-time bound.
The confidence parameter increases the required observations while
the same public interior dictionary supplies the transmitted word.
